\begin{pastexam}[label={px:ncku-110-prog-03}]{成大資訊聯招 110 學年度 第 3 題}
（計算）

(10\%) Assume that you have to store the four strings: ``abcdef'', ``bcdefa'', ``cdefab'', and
``defabc'', in the hash table by using the hashing technique.
\begin{enumerate}[label=(\arabic*),itemsep=0pt,topsep=2pt]
\item (5\%) Consider the hash function: the index for a specific string will be equal to the sum
of the ASCII values of the characters modulo 11. Please determine how many collisions if we use
this hash function.
\item (5\%) Consider the hash function: the index for a specific string will be equal to sum of
ASCII values of characters multiplied by their respective order in the string after which it is
modulo with 11. Please determine how many collisions if we use this hash function.
\end{enumerate}
Note: The ASCII values of a, b, c, d, e, and f are 97, 98, 99, 100, 101, and 102 respectively.
\end{pastexam}

\begin{pastsoln}
答案：(1) $3$ 次碰撞；(2) $0$ 次碰撞。

題目沒有說「幾次碰撞」怎麼數。 這裡照最常見的算法：依題目給的順序把字串放進表裡，放進去時
起始位置已經有別的字串，就算一次碰撞。 另一種數法是數「落在同一格的字串有幾對」，答案寫在
括號裡。

(1) 四個字串都是 a 到 f 各用一次，只是順序不同，字元碼的和都是
\[
97 + 98 + 99 + 100 + 101 + 102 = 597, \qquad 597 = 11 \times 54 + 3,
\]
四個字串的位置都是 $3$。 第一個放進去沒有碰撞，後面三個都撞上，共 $3$ 次碰撞（落在同一格
的有 $\binom{4}{2} = 6$ 對）。 這就是資結書第~12.1~節說的：只把字元碼加起來，用
同一組字母的字串一定撞在一起。

(2) 第 $i$ 個字元乘上 $i$（$i = 1, \ldots, 6$）再加起來：
\begin{center}
\small
\begin{tabular}{lll}
\toprule
字串 & 加權和 & 除以 $11$ 的餘數 \\
\midrule
abcdef & $1 \cdot 97 + 2 \cdot 98 + 3 \cdot 99 + 4 \cdot 100 + 5 \cdot 101 + 6 \cdot 102 = 2107$ & $2107 - 11 \times 191 = 6$ \\
bcdefa & $1 \cdot 98 + 2 \cdot 99 + 3 \cdot 100 + 4 \cdot 101 + 5 \cdot 102 + 6 \cdot 97 = 2092$ & $2092 - 11 \times 190 = 2$ \\
cdefab & $1 \cdot 99 + 2 \cdot 100 + 3 \cdot 101 + 4 \cdot 102 + 5 \cdot 97 + 6 \cdot 98 = 2083$ & $2083 - 11 \times 189 = 4$ \\
defabc & $1 \cdot 100 + 2 \cdot 101 + 3 \cdot 102 + 4 \cdot 97 + 5 \cdot 98 + 6 \cdot 99 = 2080$ & $2080 - 11 \times 189 = 1$ \\
\bottomrule
\end{tabular}
\end{center}
四個位置 $6, 2, 4, 1$ 都不同，$0$ 次碰撞（$0$ 對）。 乘上位置以後，字元的順序會影響結果，
所以能分開這幾個字串；資結書第~12.1~節用 $B$ 進位（Horner 法）是同樣的想法，而且
權重更分散。
\end{pastsoln}
